\subsubsection{Objetivo y metas}

\subsubsubsection{Objetivo general}
Realizar un diagnóstico integral de la infraestructura tecnológica del CIITA para identificar vulnerabilidades críticas y proponer recomendaciones técnicas que permitan modernizar los sistemas, garantizar el cumplimiento normativo y optimizar la operación del centro.

\subsubsubsection{Objetivos específicos}
\begin{itemize}
	\item \textbf{Actualización de equipos críticos}:
	\begin{itemize}
		\item Reemplazar software obsoletos en:
		\begin{itemize}
			\item 16 equipos del laboratorio de planta alta (área abierta).
			\item 24 equipos del laboratorio de planta baja (área especializada).
			\item 10 equipos de escritorio administrativos.
			\item 7 laptops asignadas a personal.
			\item 18 laptops de UMA (Unidad Móvil de Aprendizaje)
		\end{itemize}
	\end{itemize}
	
	\item \textbf{Análisis de configuración de red}:
	\begin{itemize}
		\item Identificar riesgos en switches y routers que conectan los  equipos de la institución. 
		\item Evaluar políticas de seguridad.
	\end{itemize}
	
	
	\item \textbf{Concienciación en ciberseguridad}:
	\begin{itemize}
		\item Capacitar al 100\% del personal administrativo en manejo de datos según LGPDPSO.
	\end{itemize}
	
	\item \textbf{Cumplimiento LGPDPSO}:
	\begin{itemize}
		\item Diagnosticar requisitos no cumplidos .
		\item Proponer soluciones técnicas para switches.
	\end{itemize}
\end{itemize}